\documentclass[12pt]{article}
\usepackage[margin=1in]{geometry}
\usepackage[all]{xy}


\usepackage{amsmath,amsthm,amssymb,color,latexsym}
\usepackage{geometry}        
\geometry{letterpaper}  
\usepackage{boxedminipage}
\usepackage{graphicx}
\usepackage{hyperref}
\hypersetup{
    colorlinks=true,
    citecolor=blue,
    linkcolor=blue,
    filecolor=cyan,
    urlcolor=magenta,
}
\usepackage[capitalize,nameinlink]{cleveref}

\newtheorem{problem}{Problem}
\newtheorem{lemma}{Lemma}


\newenvironment{solution}[1][\it{Solution}]{\textbf{#1. } }{$\square$}
\newcommand{\pointsbox}[1]{%
  \colorbox{red!70}{\strut\textcolor{white}{\sffamily\bfseries\ \ #1\ \ }}%
}
\newcommand{\pagetag}[1]{%
  \colorbox{red!70}{\strut\textcolor{white}{\sffamily\bfseries\ \ #1\ \ }}%
}
\newcommand{\striptitle}[2]{%
  \noindent
  \begin{minipage}{0.07\linewidth}\centering \pagetag{#1}\end{minipage}%
  \hfill
  \begin{minipage}{0.89\linewidth}
    {\large\scshape #2}
  \end{minipage}\par\vspace{0.75em}
}
% --- Environments ---
\theoremstyle{plain}
\newtheorem{theorem}{Theorem}
\newtheorem{corollary}[theorem]{Corollary}
\theoremstyle{definition}
\newtheorem{definition}[theorem]{Definition}
\newtheorem{remark}[theorem]{Remark}

\newenvironment{boxedalgo}
  {\begin{center}\begin{boxedminipage}{1\linewidth}}
  {\end{boxedminipage}\end{center}}

\newenvironment{inlinealgo}[1]
  {\noindent\hrulefill\\\noindent{\bf {#1}}}
  {\hrulefill}


% --- Shortcuts ---
\newcommand{\R}{\mathbb{R}}
\newcommand{\inner}[2]{\left\langle #1,#2 \right\rangle}
\newcommand{\proj}{\mathrm{proj}}
\newcommand{\Id}{\mathbf{I}}
\newcommand{\ot}{\otimes}
\newcommand{\mc}{\mathcal}

\begin{document}
\noindent CS 578 Statistical Machine Learning (Fall 2026)\hfill Problem Set \# 4\\
\{\{Your name\}\} \hfill \textbf{Due: October 9th}

\noindent\hrulefill

\subsection*{Problem 1. $k$-Means Clustering.}
Show that the $k$-means algorithm might converge to a point that is not a local minimum. 

\emph{Hint:} Suppose that $k = 2$ and the sample points are $\{1,2,3,4\}\subset \R$. Suppose we initialize the $k$-means with the centers $\mu_1=2, \mu_2=4$; and suppose we break ties in the definition of $C_i$ by assigning $i$ to be the smallest value in $\arg \min_j \|{\bf x}-{\bf \mu}_j\|$.

\vspace{1em}
\noindent{\bfseries Solution.}\newline 
\   %%% <-- ERASE THIS LINE AND WRITE YOUR SOLUTION HERE
\newpage

\subsection*{Problem 2. Spectral Clustering.}
In lecture, we studied spectral partitioning using the \emph{Fiedler vector}, an eigenvector corresponding to the second-smallest eigenvalue of the graph Laplacian. In this problem, we analyze a basic algorithm for computing a leading eigenvector: the \emph{power method}.
\begin{enumerate}
    \item Let $A\in\R^{n\times n}$ be a symmetric positive semidefinite matrix with eigenvalues
\[
    \lambda_1>\lambda_2\geq\cdots\geq\lambda_n\geq 0
\]
and corresponding orthonormal eigenvectors $\mathbf{v}_1,\dots,\mathbf{v}_n$. Starting from any vector $\mathbf{x}^0=\sum_{i=1}^n a_i\mathbf{v}_i$ with $a_1\neq 0$, the power method performs the iterations
\[
    \mathbf{x}^{t+1}
    =\frac{A\mathbf{x}^t}{\|A\mathbf{x}^t\|_2},
    \qquad t=0,1,\dots.
\]
Prove that, for every $i\in \{2,\dots,n\}$ and every integer $t\geq 1$,
\[
    \left|
        \frac{\langle\mathbf{x}^t,\mathbf{v}_i\rangle}
             {\langle\mathbf{x}^t,\mathbf{v}_1\rangle}
    \right|
    \leq
    \left(\frac{\lambda_i}{\lambda_1}\right)^t
    \left|
        \frac{\langle\mathbf{x}^0,\mathbf{v}_i\rangle}
             {\langle\mathbf{x}^0,\mathbf{v}_1\rangle}
    \right|.
\]
Thus, $\mathbf{x}^t$ converges to the top eigenvector $\pm{\bf v}_1$.
\item Let $\alpha=a_1^2/\|\mathbf{x}^0\|_2^2$. Prove that, for every integer $t\geq 1$,
\[
    (\mathbf{x}^t)^\top A\mathbf{x}^t
    \geq \lambda_1\alpha^{1/(2t+1)}.
\]
That is, for any $\varepsilon\in(0,1)$, every integer
$t\geq \max\{1,\frac{1}{2\varepsilon}\log(1/\alpha)\}$ satisfies
\[
    (\mathbf{x}^t)^\top A\mathbf{x}^t
    \geq (1-\varepsilon)\lambda_1.
\]
Unlike convergence to the leading eigenvector, this guarantee does not depend on the spectral gap $\lambda_1-\lambda_2$.

\emph{Hint:} Apply Jensen's inequality to the concave function
$u\mapsto u^{2t/(2t+1)}$.  
\end{enumerate}

\noindent{\bfseries Solution.}\newline 
\   %%% <-- ERASE THIS LINE AND WRITE YOUR SOLUTION HERE
\newpage

\end{document}
